\section{Comparison with representative-set and policy-count results}
\label{app:comparison}
The affine coefficient representation clarifies both the connection to
prior work and the precise scope of the structural results. This section
does not claim that additive regret creates an unrelated optimization
problem: its supporting-line geometry is shared with established
representative-set methods.

\paragraph{Convex Pareto approximation.}
For a finite coefficient menu $Q\subset\mathcal L$,
\[
 \min_{z\in Q}\langle(1,p),z\rangle
 =\min_{z\in\operatorname{conv}Q}\langle(1,p),z\rangle.
\]
Thus taking a convex hull \emph{after selecting} the menu does not change
its value at any price. Allowing its $K$ selected coefficient points to
come from $\operatorname{conv}\mathcal L$, however, can change the optimum.
This is the distinction between the representative's feasibility and
convex evaluation of the selected set. In particular, use of convex
combinations in a Pareto-approximation definition does not itself distinguish
that literature from finite-menu evaluation.

Diakonikolas' detailed treatment~\cite[Chapter~5]{diakonikolas2011},
following~\cite{diakonikolas2008}, compares representatives selected from
a finite set with representatives allowed in its convex hull. Its
Section~5.2.2 explicitly recognizes the usefulness of points away from
the lower convex envelope. Theorems~5.3.3 and~5.3.5 concern multiplicative
approximation and cardinality with a weighted-optimization oracle;
Remark~5.3.6 distinguishes objective-space points from recoverable decisions.
Our separation example is a spectral realization of an established
obstruction. Our contribution concerns the all-subinterval exactness
criterion, pure spectral face recovery, and a repair count localized to
missing-slope prices, with every tied optimal preset retained in
$\mathcal G$. Merely changing multiplicative to additive error would not
by itself establish a substantive contribution.

\paragraph{Scalarization and replacement.}
Weighted-sum approximation factors and supported-solution limitations
are established in~\cite{bazgan2022,helfrich2024}. The factor-two proof
here is their elementary endpoint-regret counterpart; the sharp examples
place that loss inside the specified spectral model. Similarly,
Magnanti and Stratila~\cite[Lemma~2]{magnanti2012} replace multiplicative
piecewise-linear approximants with tangent pieces. Our refinement preserves
an additive tolerance and counts the particular prices at which a
realizable supporting line is unavailable. Rellich's theorem supplies
spectral finiteness; neither supporting-line replacement nor analytic
eigenvalue parameterization originates in this paper.

\paragraph{Robust policy counts.}
Finite-policy reductions based on convex representation appear in
\cite[Theorem~1]{hanasusanto2015} and~\cite{buchheim2017}.
Kurtz~\cite[Theorems~2 and~7]{kurtz2026} gives exactness and approximation
bounds for objectives concave in the uncertain parameter over a convex
uncertainty set, within a bounded integer-recourse setting. Here the
candidate projectors form a continuum and
$R(P,p)=\alpha(P)+pc(P)-F(p)$ is convex in price, so the required
concavity is absent in general.

Lifting to the cost vector $(1,p,-F(p))$ does not repair that mismatch:
its image is generally nonconvex, and convexifying it can change the menu
objective. For example, the two coordinate presets of Model M have zero
regret at every actual price. The lifted endpoint vectors are $(1,0,1)$
and $(1,1,0)$. At their midpoint $(1,1/2,1/2)$, both preset costs equal
$1/2$. Thus the supremum over the lifted convex hull is at least $1/2$,
although the original menu regret is zero. These policy-count results
therefore do not directly supply \eqref{eq:repair} through a minimax swap.
Conversely, our repair bound provides none of their algorithmic complexity
guarantees. Continuous hitting and covering, including the farthest-reach
characterization, already have explicit finite-adaptability treatments
\cite{cardinal2025}.
